% ===================================================================
%  TABLEAU N° 2 — Le corps humain. — La récréation. Les jeux.
%
%  Traduction, faite en 2026, en francais standard du Quebec, sur
%  l'ido de texto/io — non sur le fac-simile francais. Les regles de
%  la colonne sont posees en tete de 10-tableau-01.tex : memes cles
%  « %%K », memes renvois, pas de \nl ni de \cc, pas de \VUpage, gras
%  sur le terme seul, et pas d'espace devant « ; », « ! », « ? ».
%
%  LES SEPT CHOIX PROPRES A CE TABLEAU :
%    (36) foot-ball ............ soccer
%    (37) bicyclette ........... vélo
%    (13) corde ................ corde à danser
%    (20) jeu de tonneau ....... jeu de la grenouille
%    (12) professeur de gymnastique . professeur d'éducation physique
%    (34) lawn-tennis .......... tennis
%         état corporel ........ condition physique
%
%  « SOCCER » N'EST PAS UNE FANTAISIE, C'EST LE MOT DU PAYS. L'ido
%  ecrit « ped-balono » et le narrateur precise que c'est un jeu
%  ANGLAIS, distinct du jeu de barres : c'est donc l'association
%  football, et non le rugby ni le football canadien. Or au Quebec
%  « football » nomme le football canadien, qui est un autre sport :
%  garder « football » ici serait nommer le mauvais jeu. « Soccer »
%  est le seul mot qui designe au Quebec ce que la gravure montre.
%
%  (20) SUIT L'IDO, NON LE FAC-SIMILE. Le livret francais dit « jeu de
%  tonneau » ; l'ido dit « rano-ludo », le jeu de la GRENOUILLE, et sa
%  note (*) explique justement que les deux noms designent le meme
%  jeu — tonneau en allemand et en francais, grenouille en anglais.
%  On traduit l'ido : c'est la grenouille. La note, elle, garde les
%  deux noms, puisque c'est d'eux qu'elle parle.
%
%  DEUX PIEGES DE NOMENCLATURE ANATOMIQUE, tranches sur l'ido et non
%  sur le francais de 1926 :
%    * « kruro » est la CUISSE et « suro » le MOLLET ; l'ido enumere
%      kruro, genuo, poplito, suro — cuisse, genou, jarret, mollet —
%      dans l'ordre descendant, et c'est ce qui fixe les deux.
%    * « maxilo » est la machoire SUPERIEURE, « mandibulo »
%      l'INFERIEURE : les deux mots ido sont distincts la ou le
%      francais courant n'a qu'une « machoire ».
% ===================================================================

%%K t02-apar-1 apar
\VUsaut{4.0mm}
\VUtitre{15.0pt}{60}{TABLEAU N\textsuperscript{o} 2}
\VUsaut{2.0mm}
\VUtitre{11.4pt}{0}{Le corps humain. --- La récréation.}
\VUtitre{11.4pt}{0}{Les jeux.}

%%K t02-tit-0 sub
\VUtitre{13.2pt}{0}{Le corps humain.}
\VUsaut{2.4mm}
\VUcentre{10.2pt}{0}{\textit{(Simple leçon de sciences naturelles.)}}
\VUsaut{3.0mm}

%%K t02-tit-1 sub pos=t02-01-1
\VUpk{10.2pt}{40}{La tête.}
\VUsaut{3.0mm}

%%K t02-01-1 p
1. --- Les principales parties du corps humain sont : la \VUgras{tête}, le
\VUgras{tronc} et les \VUgras{membres}.

%%K t02-02-1 p
2. --- La tête comprend deux parties : le \VUgras{crâne}, qui contient le
\VUgras{cerveau} et que couvrent les \VUgras{cheveux}, et le
\VUgras{visage}.

%%K t02-03-1 p
3. --- Sur notre dessin, on ne voit ni le crâne ni le cerveau; mais on voit
très clairement les parties importantes du visage.

%%K t02-04-1 p
4. --- Voici d'abord le \VUgras{front}, qui annonce une puissante
intelligence, une pensée toujours en action. Les \VUgras{tempes} sont très
dégagées. L'\VUgras{œil} est vif et grand ouvert. Un \VUgras{sourcil} épais et
de longs \VUgras{cils} le protègent.

%%K t02-05-1 p
5. --- Voici encore le \VUgras{nez} et les \VUgras{narines}. Le nez nous sert
à sentir et à respirer. De chaque côté du nez, il y a les \VUgras{joues}, et
plus loin les \VUgras{oreilles}. Le bas du visage se compose de la
\VUgras{bouche} et du \VUgras{menton}.

%%K t02-06-1 p
6. --- On ne les voit pas très bien, parce que la \VUgras{moustache} et les
\VUgras{favoris} nous les cachent. Mais on sait que la bouche a les deux
\VUgras{lèvres}, la \VUgras{mâchoire supérieure} et la \VUgras{mâchoire
inférieure} avec leurs \VUgras{dents}, et la \VUgras{langue}. La bouche nous
sert à parler et à manger. La langue nous sert à goûter les aliments.

%%K t02-07-1 p
7. --- L'œil, l'oreille, le nez et la bouche sont, comme les doigts, les
organes des cinq sens : la \VUgras{vue}, l'\VUgras{ouïe}, l'\VUgras{odorat},
le \VUgras{goût}, le \VUgras{toucher}.

%%K t02-08-1 p
8. --- La tête est donc l'une des parties les plus intéressantes du corps
humain.

%%K t02-tit-2 sub
\VUsaut{4.4mm}
\VUpk{10.2pt}{40}{Le tronc.}
\VUsaut{3.0mm}

%%K t02-09-1 p
9. --- Le \VUgras{cou} joint la tête au tronc. Il comprend la \VUgras{gorge}
et la \VUgras{nuque}.

%%K t02-10-1 p
10. --- Le tronc contient aussi beaucoup d'organes essentiels. Dans la
\VUgras{poitrine} se trouvent les \VUgras{poumons}, avec lesquels nous
respirons, et le \VUgras{cœur}, qui est le centre de la vie. Quand le cœur
cesse de battre, l'être vivant meurt.

%%K t02-11-1 p
11. --- Les \VUgras{côtes} soutiennent la poitrine.

%%K t02-12-1 p
12. --- Les autres parties du tronc sont le \VUgras{flanc}, le \VUgras{dos},
les \VUgras{épaules} et l'\VUgras{abdomen} (le ventre). Nous nous couchons sur
le flanc ou sur le dos pour dormir; nous portons les charges sur nos épaules.

%%K t02-13-1 p
13. --- Dans l'abdomen se trouvent l'\VUgras{estomac} et les
\VUgras{intestins}. Ils nous servent à digérer les aliments.

%%K t02-tit-3 sub
\VUsaut{4.4mm}
\VUpk{10.2pt}{40}{Les membres.}
\VUsaut{3.0mm}

%%K t02-14-1 p
14. --- Nous avons quatre \VUgras{membres} : deux \VUgras{bras} et deux
\VUgras{jambes}. Avec les bras, nous prenons, nous tenons, nous levons et nous
ramassons les objets; avec les jambes, nous nous tenons debout, nous marchons
et nous courons.

%%K t02-15-1 p
15. --- L'\VUgras{épaule} joint le bras au corps. Un \VUgras{muscle} très
robuste, le biceps, fait bouger le bras, que le \VUgras{coude} joint à
l'\VUgras{avant-bras}. Le \VUgras{poignet} forme l'articulation de la
\VUgras{main}, qui se termine par les \VUgras{doigts} et les \VUgras{ongles}.
Sur la main, on distingue une \VUgras{veine} (et une artère) dans laquelle
coule le \VUgras{sang}.

%%K t02-16-1 p
16. --- Les parties de la jambe sont : la \VUgras{cuisse}, le \VUgras{genou}
et le \VUgras{jarret}, le \VUgras{mollet}, dont la \VUgras{chair} est couverte
par la \VUgras{peau}, la \VUgras{cheville} et le \VUgras{pied}. Les
\VUgras{os} de la jambe sont très gros et très solides. Au bout du pied, il y a
les \VUgras{orteils}. Les autres parties sont la \VUgras{plante} et le
\VUgras{talon}.

%%K t02-17-1 p
17. --- Telle est la description à peu près complète du corps humain.

%%K t02-17-2 p
\textit{Note.} --- \textit{Avec l'expression « j'ai mal à... (à la tête, etc.)
», le professeur pourra reprendre tout ce vocabulaire du point de vue de la
bonne ou de la mauvaise} \VUgras{condition physique}.

%%K t02-tit-4 sub
\VUsaut{5.0mm}
\VUpk{10.2pt}{40}{Gymnastique.}
\VUsaut{2.4mm}
\VUcentre{10.2pt}{0}{\textit{(Récit d'un des élèves.)}}
\VUsaut{3.0mm}

%%K t02-18-1 p
18. --- Pour être souples, agiles et en bonne santé, nous devons exercer nos
jambes et nos bras. C'est pourquoi nous jouons et nous faisons de la
\VUgras{gymnastique}.

%%K t02-19-1 p
19. --- Pendant la semaine, ces deux exercices ont lieu dans la cour du lycée
(voir le tableau n\textsuperscript{o} 1). Mais le jeudi et le dimanche, nous
allons sur une grande pelouse, où nous avons de l'espace pour courir et nous
amuser.

%%K t02-20-1 p
20. --- Nos maîtres et nos parents nous y accompagnent parfois; mais c'est le
\VUgras{professeur d'éducation physique}\textsuperscript{(12)} qui dirige les
exercices.

%%K t02-21-1 p
21. --- Sur la pelouse, à gauche de l'entrée, se trouvent les \VUgras{appareils
de gymnastique}\textsuperscript{(1)}; nous grimpons à
l'\VUgras{échelle}\textsuperscript{(2)}, nous nous renversons aux
\VUgras{anneaux}\textsuperscript{(3)} et au
\VUgras{trapèze}\textsuperscript{(4)}, nous nous hissons à la force des bras
jusqu'en haut de l'\VUgras{échelle de corde}\textsuperscript{(5)} ou de la
\VUgras{corde à nœuds}\textsuperscript{(6)}.

%%K t02-22-1 p
22. --- Pendant ce temps, nos camarades sautent par-dessus le \VUgras{cheval de
bois}\textsuperscript{(7)} ou les \VUgras{barres
parallèles}\textsuperscript{(11)}. D'autres \VUgras{boxent}\textsuperscript{(8)}
ou font de l'\VUgras{escrime}\textsuperscript{(9)} avec un masque et un
\VUgras{fleuret}. D'autres enfin lèvent des
\VUgras{haltères}\textsuperscript{(10)}.

%%K t02-tit-5 sub
\VUsaut{4.6mm}
\VUpk{10.2pt}{40}{Les jeux.}
\VUsaut{3.0mm}

%%K t02-23-1 p
23. --- Autour des gymnastes, des garçons et des filles jouent à différents
\VUgras{jeux}. Les filles s'amusent avec leur
\VUgras{cerceau}\textsuperscript{(14)}; elles sautent à la \VUgras{corde à
danser}\textsuperscript{(13)} ou invitent leurs frères et leurs cousins à une
partie de \VUgras{croquet}\textsuperscript{(15)}. Elles adorent chasser la
\VUgras{boule} de l'adversaire avec leur \VUgras{maillet de bois}, juste au
moment où il croit passer facilement sous l'arceau!

%%K t02-24-1 p
24. --- Les garçons préfèrent les jeux plus physiques. Ils jouent volontiers
aux \VUgras{quilles}\textsuperscript{(16)}, qu'ils abattent adroitement avec
une \VUgras{boule}\textsuperscript{(17)}. Ils ne dédaignent pas non plus la
\VUgras{toupie}\textsuperscript{(18)} et le
\VUgras{bilboquet}\textsuperscript{(19)}, ni même le \VUgras{jeu de la
grenouille}\textsuperscript{(20)}. Mais leurs jeux préférés sont les
\VUgras{billes}\textsuperscript{(21)} et la \VUgras{balle au
mur}\textsuperscript{(22)}, parce que le mur de la
\VUgras{serre}\textsuperscript{(26)} est très commode pour faire rebondir la
balle légère.

%%K t02-25-1 p
25. --- Le petit Jean, monté sur ses \VUgras{échasses}\textsuperscript{(24)},
est très adroit et sait très bien garder l'équilibre. Près de lui, quelques
camarades jouent à \VUgras{cache-cache}\textsuperscript{(25)} dans la serre et
derrière la \VUgras{haie}. D'autres préfèrent la \VUgras{main
chaude}\textsuperscript{(28)} ou le \VUgras{volant}\textsuperscript{(29)}, qui
voltige d'une \VUgras{raquette} à l'autre.

%%K t02-noto-1 noto
\VUnotes{143.2mm}{%
(*) Les jeux d'enfants ont souvent des noms purement nationaux. Nous avons
cherché à les nommer en ido du mieux que nous pouvions, soit en nous inspirant
de l'action même qui les caractérise, soit en choisissant le nom national de
tel ou tel pays, quand il n'est pas trop impropre. Par exemple : le \VUgras{jeu
du tonneau} (en allemand et en français) est le \VUgras{jeu de la grenouille}
\textsuperscript{(20)} (en anglais), où les joueurs s'efforcent de lancer leurs
disques ronds dans la bouche de la grenouille de métal qui se tient bouche bée
sur la caisse (le tonneau) percée de trous.
Le \VUgras{cheval fondu}\textsuperscript{(30)} ne diffère du
\VUgras{saute-mouton} que par ceci : pour sauter par-dessus la tête, les
joueurs appuient les mains sur les épaules de leur partenaire, tandis que dans
le
« \VUgras{saute-mouton} » ils n'appuient les mains que sur son dos. --- Ou
encore, quelques-uns des participants se penchent pendant que les autres
sauteront par-dessus leur dos en appuyant les mains sur l'épaule; les premiers
doivent supporter le choc sans fléchir, les seconds doivent sauter sans perdre
l'équilibre (voir le tableau lui-même).%
}

%%K t02-26-1 p
26. --- Au milieu de la pelouse, il y a deux arbres. Au pied de l'un des deux,
sept ou huit garçons ont organisé une partie de \VUgras{cheval
fondu}\textsuperscript{(30)} (*). Ce jeu n'est pas connu dans tous les pays;
mais tout le monde sait ce qu'est le \VUgras{saute-mouton}, auquel les enfants
ne jouent pas cette fois-ci.

%%K t02-tit-6 sub
\VUsaut{5.0mm}
\VUpk{10.2pt}{40}{Les jeux en plein air.}
\VUsaut{3.4mm}

%%K t02-27-1 p
27. Le \VUgras{colin-maillard}\textsuperscript{(31)} est aussi très amusant;
filles et garçons mettent tour à tour le \VUgras{bandeau} sur leurs yeux et
attrapent les maladroits qui se laissent prendre.

%%K t02-28-1 p
28. --- Au milieu de la pelouse, voici un jeu bien français, mais un peu
violent : c'est l'\VUgras{ourse}\textsuperscript{(32)}, beaucoup plus bruyant
que le jeu si tranquille du \VUgras{cerf-volant}\textsuperscript{(33)}, ou même
que le jeu anglais du \VUgras{tennis}\textsuperscript{(34)}.

%%K t02-29-1 p
29. --- Ce dernier sport fait la joie des grands élèves. Nos professeurs
eux-mêmes le pratiquent volontiers. Il demande beaucoup d'adresse et
d'agilité, et il me plaît beaucoup. Je laisse aux filles les plaisirs de la
\VUgras{balançoire}\textsuperscript{(35)}.

%%K t02-30-1 p
30. --- Je n'aime pas un autre jeu anglais, qui pourrait devenir dangereux.

%%K t02-30-1b p
Je veux parler du \VUgras{soccer}\textsuperscript{(36)}, qui fait le bonheur de
mon frère aîné. Je préfère notre vieux \VUgras{jeu de
barres}\textsuperscript{(38)}, tout aussi bon pour la santé et vraiment moins
brutal.

%%K t02-tit-7 sub
\VUsaut{4.6mm}
\VUpk{10.2pt}{40}{Vélo et jeux de ballon.}
\VUsaut{3.0mm}

%%K t02-31-1 p
31. --- J'aime aussi faire le tour de la pelouse à
\VUgras{vélo}\textsuperscript{(37)}.

%%K t02-32-1 p
32. --- L'année prochaine, j'apprendrai
l'\VUgras{équitation}\textsuperscript{(39)}. Qu'il est beau, un cheval qui
marche ou qui trotte avec grâce, et qui saute les obstacles au galop!

%%K t02-33-1 p
33. --- L'été, nous apprenons la \VUgras{natation}\textsuperscript{(40)}. Nous
plongeons et nous nous baignons avec délices dans l'eau claire et fraîche de la
rivière. Parfois, enfin, nous lançons un \VUgras{ballon} de papier
\textsuperscript{(41)} qui monte majestueusement dans le ciel dès qu'il est
rempli d'air chaud.

%%K t02-34-1 p
34. --- Il y a encore bien d'autres jeux, mais je n'en finirais pas de les
énumérer, et l'on comprend, d'après ce résumé, qu'on ne s'ennuie pas beaucoup,
le jeudi, sur notre belle pelouse.

%%K t02-34-2 p
N.B. --- \textit{Le professeur complétera facilement ce chapitre en décrivant
plus en détail chacun des jeux les plus importants.}
\VUsaut{6.0mm}
\VUfilet{22mm}
